\begin{table}[t]
\centering
\caption{Boost-converter parameters (averaged and switching-level models).}
\label{tab:Table_Converter_Parameters}
\small
\begin{tabular}{lll}
\hline
parameter & value & note \\
\hline
L & 0.001 & chosen so that the ripple at the nominal operating point (V\_pv = 26 V, d = 0.46, f\_sw = 20 kHz) is dI = V\_pv d / (f\_sw L) = 0.6 A (about 7 \% of 9 A); CCM holds above 0.3 A \\
R\_L & 0.02 & inductor series resistance \\
C\_in & 0.0005 & input ripple dV = dI / (8 f\_sw C\_in) = 0.08 V; also fixes the input-capacitor time constant of the voltage loop \\
f\_sw & 20000 & typical for a 250 W module-level converter \\
V\_bus & 48 & constant 48 V bus (battery or regulated DC link), the same value used for the display-only duty of the quasi-static results \\
V\_f & 0.5 & diode forward drop \\
R\_on & 0.015 & switch on-resistance \\
Kp & 0.05 & PI proportional gain [1/V], fixed once on the step response at G = 800 W/m2, T = 25 C (tools in this module), before any scenario was simulated \\
Ki & 10 & PI integral gain [1/(V s)] \\
D\_MIN & 0.05 &  \\
D\_MAX & 0.95 &  \\
W\_MS & 20 & settling window per measurement [ms] (3x for command jumps larger than 2 V); one MPPT control cycle (1 s) holds at least one window; a 14-point scan needs at most 14 x 60 ms = 0.84 s \\
DT\_AVG & 2e-05 & integration step of the averaged model [s] \\
STEPS\_PER\_PERIOD & 100 & switching-level steps per PWM period \\
\hline
\end{tabular}
\end{table}
